\documentclass[10pt,a4paper]{amsart}
\usepackage[margin=19mm]{geometry}
\usepackage{amsmath,amssymb,mathtools}
\usepackage[scr=rsfs,cal=euler]{mathalfa}
\usepackage{xfrac}
\usepackage[unicode,hidelinks,bookmarks=false]{hyperref}
\newcommand{\ie}{i.e.\ }
\newcommand{\cf}{cf.\ }
\newcommand{\resp}{respectively\ }
\newcommand{\Subsection}{\subsection}
\providecommand{\nobreakdash}{\nobreak\hbox{-}}
\setlength{\emergencystretch}{2em}
\multlinegap0pt
\usepackage{bm}
\usepackage{bbm}
\usepackage{dsfont}
\usepackage{xspace}
\usepackage{cancel}
\usepackage{mathtools}
\usepackage{cleveref}
\usepackage{amsthm}
\makeatletter
\@ifundefined{theorem}{\newtheorem{theorem}{Theorem}}{}
\@ifundefined{lemma}{\newtheorem{lemma}[theorem]{Lemma}}{}
\makeatother
\DeclareMathOperator{\tr}{tr}
\def\bigopen#1{\left(#1\right)}
\def\mA{{\bm{A}}}
\def\mC{{\bm{C}}}
\def\mI{{\bm{I}}}
\def\mT{{\bm{T}}}
\def\mX{{\bm{X}}}
\def\vone{{\bm{1}}}
\def\vw{{\bm{w}}}
\def\vzero{{\bm{0}}}
\newcommand{\E}{\mathbb{E}}
\newcommand{\MARPCD}{\mathcal{M}_{\bm{A}}^{\RPCD}} % JW: Use emph in theorems
\newcommand{\RPCD}{\text{RPCD}} % JW: Use emph in theorems
\title[Provable Benefit of Random Permutations over Uniform Samplin]{Provable Benefit of Random Permutations over Uniform Sampling in Stochastic Coordinate Descent}
\author{Donghwa Kim, Jaewook Lee, Chulhee Yun}
\date{Source: arXiv:2505.23152v1 (CC BY 4.0)}
\begin{document}
\maketitle

\noindent\emph{Source and licence.} Provable Benefit of Random Permutations over Uniform Sampling in Stochastic Coordinate Descent. Version arXiv:2505.23152v1 (CC BY 4.0), distributed under the Creative Commons Attribution 4.0 International licence, as recorded in the arXiv licence metadata for this exact version and shown on the abstract page https://arxiv.org/abs/2505.23152v1. Full rights notice: licenses/arxiv-2505.23152-CC-BY-4.0.txt.\par\smallskip

\noindent\textit{Editorial scope.} Counted unit: the Lemma whose source label is \texttt{lem:taunonnegative} in the article (source0.tex lines 1763--1766, source label \texttt{lem:taunonnegative}), delivered together with the article's own complete original proof of it (source0.tex lines 1767--1870, one \texttt{proof} environment of 104 lines). No mathematics is written, repaired or re-derived: statement and proof are the article's own text line for line. Required context retained uncounted: none. Every definition, display and label of those lines is retained unchanged. Omitted from this package and not redistributed: 1--1762, 1871--4293. Nothing inside a retained excerpt is omitted or altered: every retained body is delivered exactly as the article's lines are written, so no omission receipt of any kind accompanies this package. Citations: the retained excerpts carry no citation command at all, so no bibliography entry of the article is transcribed and no reference is redistributed. Reference adaptation: none was needed and none was made; every reference inside the delivered unit resolves inside it, so no unresolved reference or citation is produced. Printed slips kept literal: no printed word, symbol, hypothesis or number of the counted statement or of its proof was altered. Layout and class: the article's own class is \texttt{article} with packages including microtype, graphicx, subfigure, booktabs, xcolor, mdframed, subcaption, enumitem, hyperref, icml2025, amsmath, amssymb, mathtools, amsthm; the wrapper is the lane's pinned \texttt{amsart} editorial wrapper at 10pt on A4, which restates the theorem-like environments the retained text needs and adds the article's own macro definitions that the retained text needs (\texttt{\textbackslash E}, \texttt{\textbackslash MARPCD}, \texttt{\textbackslash RPCD}, \texttt{\textbackslash bigopen}, \texttt{\textbackslash mA}, \texttt{\textbackslash mC}, \texttt{\textbackslash mI}, \texttt{\textbackslash mT}, \texttt{\textbackslash mX}, \texttt{\textbackslash tr}, \texttt{\textbackslash vone}, \texttt{\textbackslash vw}, \texttt{\textbackslash vzero}), with extra wrapper packages (bm, bbm, dsfont, xspace, cancel, mathtools, cleveref). The delivered numbering is the unit's own. Assets: no retained span contains a figure environment, a graphics-inclusion command, a PDF-page inclusion, a table or a TikZ picture; no image file is copied into this package, the package declares no asset dependency and no image file is redistributed with it. Licence: version arXiv:2505.23152v1 of this article is distributed under the Creative Commons Attribution 4.0 International licence, as recorded in the arXiv licence metadata for this exact version and shown on the abstract page https://arxiv.org/abs/2505.23152v1. A full rights notice is delivered at licenses/arxiv-2505.23152-CC-BY-4.0.txt. Characters: every retained excerpt is pure ASCII, so the delivered engine, XeLaTeX with its default Latin Modern text font, reports no missing character.
\begin{lemma}
    \label{lem:taunonnegative}
    $\tau_1^{(1)}, \tau_1^{(2)}, \tau_2^{(1)}$ and $\tau_2^{(2)}$ are all non-negative.
\end{lemma}

\begin{proof}
    First, we show that $\tau_1^{(1)}, \tau_1^{(2)} \ge 0$ from the fact that $\MARPCD$ preserves positive semi-definiteness. 
    
    Recall that $\MARPCD(\mX) = \E[\mT_{\mA, p}^{\RPCD \top} \mX \mT_{\mA, p}^{\RPCD}]$.
    Since $\mX \succeq 0$ implies $\mT_{\mA, p}^{\RPCD \top} \mX \mT_{\mA, p}^{\RPCD} \succeq 0$ and the expectation of positive semi-definite matrices is also positive semi-definite, we can conclude that $\MARPCD(\mX) \succeq 0$ whenever $\mX \succeq 0$.
    As both $\mI \succeq 0$ and $\vone \vone^{\top} \succeq 0$, it follows that $\MARPCD(\mI) \succeq 0$ and $\MARPCD(\vone \vone^{\top}) \succeq 0$, i.e., their eigenvalues are non-negative.
    Since $\tau_1^{(1)}$ and $\tau_1^{(2)}$ are eigenvalues of $\MARPCD(\mI) \succeq 0$ and $\MARPCD(\vone \vone^{\top}) \succeq 0$ respectively, we have $\tau_1^{(1)}, \tau_1^{(2)} \ge 0$.
    
    Next, we show that $\tau_2^{(2)}\ge 0$. For $\vw = \mC^{\top} \vone$, we have 
    \begin{align*}
        \gamma &= \vone^{\top} \vw \vw^{\top} \vone = (\vone^{\top} \vw)^2 \\
        \delta &= \tr(\vw \vw^{\top}) = \tr(\vw^{\top} \vw) = \|\vw\|^2.
    \end{align*}
    Since $\tau_2^{(2)} = \gamma - \delta$, we have
    \begin{align*}
        \tau_2^{(2)} &= (\vone^{\top} \vw)^2 - \|\vw\|^2 \\
        &= 2\sum_{i<j} w(i) w(j).
    \end{align*}
    Since $w(i) = \sigma^n - \sigma^{n+1-i} \le 0$ for all $i\in [n]$, we have $w(i) w(j) \ge 0$, and therefore $\tau_2^{(2)} = 2\sum_{i<j} w(i) w(j) \ge 0$.

    Finally, to show that $\tau_2^{(1)}\ge 0$, we first show that the entries of $\mC^{\top} \mC$ is non-negative.
    
    In particular, we can prove the following two inequalities:
    % \begin{align*}
    %     C_i^{\top} C_j \ge C_i^{\top} C_{j+1} & \quad \text{for } 2 \le i < j < n, \\
    %     C_i^{\top} C_j \ge C_{i-1}^{\top} C_j & \quad \text{for } 2 < i < j \le n,
    % \end{align*}
    \begin{itemize}
        \item If $2 \le i < j < n$, then $\mC_i^{\top} \mC_j \ge \mC_i^{\top} \mC_{j+1}$,
        \item If $2 < i < j \le n$, then $\mC_i^{\top} \mC_j \ge \mC_{i-1}^{\top} \mC_j$,
    \end{itemize}
    where $\mC_i$ denotes the $i$-th column of $\mC$.
    
    To show the first inequality, we begin by considering the difference $\mC_j - \mC_{j+1}$.
    We have
    \begin{align}
        (\mC_j - \mC_{j+1})_k &= 
        \begin{cases}
            0 &\quad \text{if } k < j, \\
            1-\sigma &\quad \text{if } k = j, \\
            (1-\sigma)(\sigma^{k-j} - \sigma^{k-j-1}) &\quad \text{if } k > j.
        \end{cases}
    \label{eq:cjdiff}
    \end{align}
    Therefore, if $2 \le i < j < n$, 
    \begin{align*}
        \mC_i^{\top} (\mC_j - \mC_{j+1}) &= (1-\sigma)^2(\sigma^{j-i} - \sigma^{j-1}) + \sum_{k=j+1}^n (1-\sigma)^2(\sigma^{k-i} - \sigma^{k-1})(\sigma^{k-j} - \sigma^{k-j-1})\\
        &= (1-\sigma)^2(\sigma^{-i} - \sigma^{-1}) \bigopen{\sigma^j +(\sigma^{-j} - \sigma^{-j-1})\sum_{k=j+1}^n \sigma^{2k}}.
    \end{align*}
    Since $i \ge 2$ and $\sigma\in (0, 1]$, we have $(1-\sigma)^2 \ge 0$ and $(\sigma^{-i}-\sigma^{-1}) \ge 0$. 
    Thus, it suffices to show that 
    \begin{align*}
        \sigma^j +(\sigma^{-j} - \sigma^{-j-1})\sum_{k=j+1}^n \sigma^{2k} &\ge 0.
    \end{align*}
    Since 
    \begin{align*}
        \sigma^j +(\sigma^{-j} - \sigma^{-j-1})\sum_{k=j+1}^n \sigma^{2k} &= \sigma^j +(\sigma^{-j} - \sigma^{-j-1}) \sigma^{2j+2} \cdot \frac{1-\sigma^{2(n-j)}}{1-\sigma^2} \\
        &= \sigma^j \bigopen{1 - \frac{\sigma}{1+\sigma}(1-\sigma^{2(n-j)})}\\
        &= \frac{\sigma^j}{1+\sigma} (1+\sigma^{2n-2j+1}),
    \end{align*}
    we have $\mC_i^{\top} (\mC_j - \mC_{j+1}) \ge 0$.

    For the second inequality, we will show that $(\mC_{i-1}-\mC_i)^{\top} \mC_j \le 0$. Using \cref{eq:cjdiff}, we obtain
    \begin{align*}
        (\mC_{i-1}-\mC_i)^{\top} \mC_j &= (1-\sigma)^2 \bigopen{-\sigma^{i-2} + (\sigma^{-i-1}-\sigma^{-i})\sum_{k=i}^{j-1} \sigma^{2k} + (\sigma^{-i+1}-\sigma^{-i})(\sigma^{-j}-\sigma^{-1})\sum_{k=j}^n \sigma^{2k}}.
    \end{align*}
    Since 
    \begin{align*}
         (\sigma^{-i+1}-\sigma^{-i})(\sigma^{-j}-\sigma^{-1})\sum_{k=j}^n \sigma^{2k} &\le 0,
    \end{align*}
    it is sufficient to show that 
    \begin{align*}
        -\sigma^{i-2} + (\sigma^{-i-1}-\sigma^{-i})\sum_{k=i}^{j-1} \sigma^{2k} \le 0.
    \end{align*}
    Since
    \begin{align*}
        -\sigma^{i-2} + (\sigma^{-i-1}-\sigma^{-i})\sum_{k=i}^{j-1} \sigma^{2k} &= -\sigma^{i-2} + (\sigma^{-i-1}-\sigma^{-i}) \sigma^{2i} \cdot \frac{1-\sigma^{2(j-i)}}{1-\sigma^2} \\
        &\le -\sigma^{i-2} + \frac{\sigma^{i-1}-\sigma^i}{1-\sigma^2}\\
        &= -\frac{\sigma^{i-2}}{\sigma+1} \\
        &\le 0,
    \end{align*}
    it follows that $(\mC_i-\mC_{i-1})^{\top} \mC_j \le 0$.

    Using a chain of the inequalities we showed, we can obtain $\min_{2\le i<j\le n} \mC_i^{\top} \mC_j = \mC_2^{\top} \mC_n$.
    Moreover, since 
    \begin{align*}
        \mC_2^{\top} \mC_n &= (1-\sigma)^2 \bigopen{1 - (\sigma^{-3}-\sigma^{-2})\sum_{k=2}^{n-1}\sigma^{2k} + (1-\sigma^{n-1})(\sigma^{n-2}-\sigma^{n-1})}
    \end{align*}
    and 
    \begin{align*}
        1 - (\sigma^{-3}-\sigma^{-2})\sum_{k=2}^{n-1}\sigma^{2k} + (1-\sigma^{n-1})(\sigma^{n-2}-\sigma^{n-1}) &\ge 1 - (\sigma^{-3}-\sigma^{-2})\sum_{k=2}^{n-1}\sigma^{2k}\\
        &\ge 1 - (\sigma^{-3}-\sigma^{-2})\frac{\sigma^4}{1-\sigma^2}\\
        &\ge 0,
    \end{align*}
    we have $\mC_2^{\top} \mC_n \ge 0$. Additionally, since $\mC_1 = \vzero$, we have $\mC_i^{\top} \mC_j = 0$ whenever $i = 1$ or $j = 1$.
    With the fact that $\mC_i^{\top} \mC_j \ge 0$ if $2 \le i, j \le n$, we have $\min_{1\le i,j \le n} \mC_i^{\top} \mC_j = \min_{1\le i,j \le n} (\mC^{\top} \mC)_{ij} \ge 0$, i.e., all entries of $\mC^{\top} \mC$ are non-negative.
    Therefore we can conclude that
    \begin{align*}
        \tau_2^{(1)} &= \vone^{\top} \mC^{\top} \mC \vone - \tr (\mC^{\top} \mC)\\
        &= \sum_{i\ne j} (\mC^{\top} \mC)_{ij}\\
        &\ge 0,
    \end{align*}
    which completes the proof.
\end{proof}

\end{document}
